% ====================================================================
% Proyecto Final de Estadística - Puntos 5, 6 y 7
% Fórmulas usadas EXCLUSIVAMENTE de:
%   [F] "Formulario de Probabilidad y Estadística" (J. D. Fuentes Mora)
%   [T] "Tabla 4: Áreas de curva normal" (R. E. Walpole) - P(Z > z)
% Cálculos realizados con R (R/calculos/parametros.R).
% Compilar con: tectonic informe.tex  (o pdflatex informe.tex)
% ====================================================================
\documentclass[11pt,a4paper]{article}

\usepackage[spanish]{babel}
\usepackage[utf8]{inputenc}
\usepackage[T1]{fontenc}
\usepackage{amsmath,amssymb,amsthm}
\usepackage{geometry}
\usepackage{booktabs}
\usepackage{array}
\usepackage{makecell}
\usepackage{xcolor}

\geometry{margin=2.5cm}
\pagestyle{plain}

% Todas las cifras numéricas vienen de R (parametros.R):
% ---- Punto 5 (fórmulas del Formulario Fuentes Mora) ----
\def\FilasTablaAgr{$[2.40,\,3.46)$ & $2.93$ & $2$ & $0.0333$ & $0.0333$ \\ $[3.46,\,4.51)$ & $3.99$ & $6$ & $0.1000$ & $0.1333$ \\ $[4.51,\,5.57)$ & $5.04$ & $12$ & $0.2000$ & $0.3333$ \\ $[5.57,\,6.63)$ & $6.10$ & $11$ & $0.1833$ & $0.5167$ \\ $[6.63,\,7.69)$ & $\mathbf{7.16}$ & $13$ & $0.2167$ & $0.7333$ \\ $[7.69,\,8.74)$ & $8.21$ & $10$ & $0.1667$ & $0.9000$ \\ $[8.74,\,9.80]$ & $9.27$ & $6$ & $0.1000$ & $1.0000$ \\}
\def\nMuestra{60}
\def\mediaM{6.5433}
\def\medianaM{6.5000}
\def\xDatosMin{2.4}
\def\xDatosMax{9.8}
\def\RTotal{7.4000}
\def\kClases{7}
\def\AClases{1.0571}
\def\mediaAgrup{6.4700}
\def\modaM{7.0514}
\def\modaLinf{6.63}
\def\modaLsup{7.69}
\def\modaFmod{13}
\def\modaFant{11}
\def\modaFpos{10}
\def\modaDuno{2}
\def\modaDdos{3}
\def\varM{2.9008}
\def\sdM{1.7032}
\def\cvM{26.03}
\def\sesgoM{-0.1231}
\def\curtM{2.5292}
\def\curtExcesoM{-0.4708}
% ---- Punto 6 (fórmulas del Formulario Fuentes Mora) ----
\def\nPob{8673}
\def\mediaP{6.7091}
\def\varP{2.9828}
\def\sdP{1.7271}
\def\cvP{25.74}
\def\errAbsMedia{0.1658}
\def\errRelMedia{2.47}
\def\errAbsVar{0.0820}
\def\errRelVar{2.75}
\def\errAbsSD{0.0239}
\def\errRelSD{1.38}
\def\errAbsCV{0.29}
\def\errRelCV{1.11}
\def\fracMuestral{0.0069}
\def\correccionFPC{0.9932}
\def\varXbarMedia{0.0497}
\def\seXbarMedia{0.2230}
\def\seVsObs{0.74}
\newcommand{\CondicionFPC}{\textbf{no} se cumple la condición del
  formulario ($n/N = \fracMuestral \leq 0{,}05$), por lo que la
  corrección por población finita \textbf{no aplica} y se usa
  directamente la fórmula de distribuciones muestrales}
\newcommand{\FormulaVarXbar}{\operatorname{Var}(\bar{X}) =
  \frac{\sigma^{2}}{\nMuestra} = \frac{\varP}{\nMuestra} = \varXbarMedia}
% ---- Punto 7 - Exp. 1: Binomial ----
\def\expPbin{0.1064}
\def\binPdos{0.2072}
\def\binPgeuno{0.6754}
\def\binE{1.0642}
\def\binV{0.9510}
% ---- Punto 7 - Exp. 2: Poisson ----
\def\poiLam{10.1429}
\def\poiPcero{0.0000}
\def\poiPgeDos{0.9996}
% ---- Punto 7 - Exp. 3: Geométrica ----
\def\geoPseis{0.0606}
\def\geoPlecuatro{0.3624}
\def\geoE{9.3965}
\def\geoV{78.8983}
% ---- Punto 7 - Exp. 4: Hipergeométrica ----
\def\hipN{8673}
\def\hipK{6287}
\def\hipn{20}
\def\hipPgeDieciocho{0.0575}
\def\hipPveinte{0.0016}
\def\hipE{14.4979}
\def\hipV{3.9797}
% ---- Punto 7 - Exp. 5: Normal (valores de la Tabla 4 de Walpole) ----
\def\normZseis{-0.32}
\def\normZseisAbs{0.32}
\def\normZocho{0.86}
\def\normZnueve{1.44}
\def\normZcinco{-0.91}
\def\normZcincoAbs{0.91}
\def\normColaZseis{0.3745}
\def\normColaZocho{0.1949}
\def\normColaZnueve{0.0749}
\def\normColaZcinco{0.1814}
\def\normPseisOcho{0.4306}
\def\probNormMayorNueve{0.0749}
\def\probNormMenorCinco{0.1814}


\newcommand{\FP}{\text{[F]}}
\newcommand{\TP}{\text{[T]}}

\title{\textbf{Proyecto Final de Estadística}\\[4pt]
\large Puntos 5, 6 y 7: parámetros descriptivos, parámetros\\
poblacionales y distribuciones de probabilidad}
\author{Luis Mango Ataulfo}
\date{Análisis de vulnerabilidades CVE (NIST NVD, enero--abril 2024)}

\begin{document}
\maketitle

\begin{abstract}
\noindent Se reportan los puntos 5, 6 y 7 del proyecto sobre la variable
$X$ = \emph{Base Score} (puntaje CVSS v3.1, escala 0.0--10.0) de
vulnerabilidades publicadas por el NIST entre enero y abril de 2024. Se trabaja
con una muestra de $\nMuestra$ observaciones extraída por \textbf{muestreo
aleatorio simple sin reemplazo} de la población de $\nPob$ registros
(\texttt{set.seed(42)}, reproducible). \textbf{Todas las fórmulas empleadas provienen exclusivamente del
formulario del curso \FP} (J.~D.~Fuentes Mora) y de la \textbf{Tabla 4 de
áreas de curva normal \TP} (R.~E.~Walpole); dicha tabla entrega
$P(Z>z)$ (área en cola derecha) y, \emph{``para valores negativos de $z$, las
áreas se encuentran por simetría''}. Todos los cálculos numéricos fueron
realizados con \texttt{R} en \texttt{R/calculos/parametros.R}, el cual exporta
las cifras que este documento importa automáticamente.
\end{abstract}

\tableofcontents
\newpage

% ====================================================================
\section{Punto 5. Parámetros descriptivos de la muestra}
% ====================================================================

Sea $\{x_1,\dots,x_{\nMuestra}\}$ la muestra de tamaño $\nMuestra$ de la
variable $X$ (Base Score). Se usan las fórmulas del formulario \FP:

\subsection{Medidas de tendencia central (datos no agrupados)}

\[
\bar{x} = \frac{1}{\nMuestra}\sum_{i=1}^{\nMuestra} x_i
\qquad\qquad
\text{Med} = \frac{x_{(n/2)} + x_{(n/2+1)}}{2}\quad (n\ \text{par})
\]

\subsection{Amplitud de intervalos (datos continuos)}

Para agrupar los datos se usan las fórmulas del formulario \FP:
\[
R = x_{\max}-x_{\min},
\qquad
k = 1 + 3{,}322\log_{10}(n),
\qquad
A = \frac{R}{k},
\qquad
m_i = \frac{L_i+L_s}{2}
\]
Con la muestra: $x_{\min}=\xDatosMin$, $x_{\max}=\xDatosMax$, por lo tanto
\[
R = \xDatosMax-\xDatosMin = \RTotal,
\qquad
k = 1+3{,}322\log_{10}(\nMuestra) = \kClases\ \text{clases},
\qquad
A = \frac{\RTotal}{\kClases} = \AClases
\]

\subsection{Moda (datos agrupados)}

Aplicando la fórmula del formulario sobre los $\kClases$ intervalos de
amplitud $A=\AClases$:
\[
\text{Mod} = L_{\inf} + \frac{d_1}{d_1+d_2}\cdot A,
\qquad d_1 = f_{\text{mod}} - f_{\text{ant}},
\qquad d_2 = f_{\text{mod}} - f_{\text{pos}}
\]

\subsection{Medidas de dispersión (datos no agrupados)}

\[
s^{2} = \frac{1}{\nMuestra-1}\sum_{i=1}^{\nMuestra}(x_i-\bar{x})^{2},
\qquad
s = \sqrt{s^{2}},
\qquad
CV = \frac{s}{\bar{x}}\times 100\,\text{\%}
\]

\subsection{Momentos y coeficientes de forma}

El formulario define los momentos; sobre los momentos centrales de la muestra
se construyen los coeficientes de sesgo y curtosis:
\[
m_k = \frac{1}{\nMuestra}\sum_{i=1}^{\nMuestra}(x_i-\bar{x})^{k}
\qquad
\text{Sesgo} = \frac{m_3}{m_2^{3/2}},
\qquad
\text{Curtosis} = \frac{m_4}{m_2^{2}},
\qquad
\text{Exceso} = \text{Curtosis} - 3
\]

\subsection{Resultados (calculados en R)}

\begin{table}[h!]
\centering
\caption{Parámetros descriptivos de la muestra de Base Score
($\nMuestra$ vulnerabilidades). Fórmulas del formulario \FP.}
\medskip
\begin{tabular}{@{}lc@{\qquad}lc@{}}
\toprule
\textbf{Estadístico} & \textbf{Valor} &
\textbf{Estadístico} & \textbf{Valor} \\
\midrule
Media ($\bar{x}$)                & $\mediaM$ &
Varianza muestral ($s^{2}$)      & $\varM$ \\
Mediana                          & $\medianaM$ &
Desv.\ estándar ($s=\sqrt{s^2}$) & $\sdM$ \\
Moda (agrupados)                 & $\modaM$ &
$CV = \frac{s}{\bar{x}}\times 100\,\text{\%}$ & $\cvM\,\text{\%}$ \\
\midrule
Sesgo $\left(\frac{m_3}{m_2^{3/2}}\right)$ & $\sesgoM$ &
Curtosis $\left(\frac{m_4}{m_2^{2}}\right)$ & $\curtM$ \\
Exceso de curtosis               & $\curtExcesoM$ & & \\
\bottomrule
\end{tabular}
\end{table}

\noindent Detalle de la moda (fórmula agrupada): el intervalo modal es
$[\modaLinf,\,\modaLsup)$ con $f_{\text{mod}}=\modaFmod$,
$f_{\text{ant}}=\modaFant$, $f_{\text{pos}}=\modaFpos$,
de donde $d_1 = \modaFmod-\modaFant = \modaDuno$,
$d_2 = \modaFmod-\modaFpos = \modaDdos$ y
\[
\text{Mod} = \modaLinf +
\frac{\modaDuno}{\modaDuno+\modaDdos}\times \AClases = \modaM
\]

\begin{table}[h!]
\centering
\caption{Tabla de frecuencias agrupadas (intervalos de amplitud
$A=\AClases$, marcas de clase $m_i$). Intervalo modal en negritas.}
\medskip
\begin{tabular}{@{}lcccc@{}}
\toprule
\textbf{Intervalo} & \textbf{Marca} & \textbf{$f_i$} &
\textbf{$f_i/n$} & \textbf{$F_i/n$} \\
\midrule
\FilasTablaAgr
\bottomrule
\end{tabular}
\end{table}

\noindent Control: la media agrupada $\bar{x}=\sum f_i m_i/n =
\mediaAgrup$ se aproxima a la media de los datos sin agrupar ($\mediaM$).

\subsection{Interpretación}

\begin{itemize}
  \item \textbf{Tendencia central:} la media ($\mediaM$) y la mediana
        ($\medianaM$) son muy próximas, lo que indica una distribución casi
        simétrica alrededor de un puntaje cercano a $6.5$, de severidad
        \emph{MEDIA}--\emph{HIGH}. La moda agrupada ($\modaM$) es inferior
        porque depende del intervalo modal $[\modaLinf,\,\modaLsup)$ y de sus
        frecuencias vecinas.
  \item \textbf{Dispersión:} la desviación estándar $s=\sdM$ y un
        $CV=\cvM\,\text{\%}$ indican dispersión moderada respecto a la media.
  \item \textbf{Forma:} el sesgo $m_3/m_2^{3/2}=\sesgoM$ es cercano a cero y
        ligeramente negativo: distribución casi simétrica con mínimo
        alargamiento de la cola izquierda. La curtosis $\curtM<3$ (exceso
        $\curtExcesoM<0$) indica una distribución \textbf{platicúrtica}: colas
        más ligeras y pico más aplanado que la normal.
  \item \textbf{Nota sobre el tamaño muestral:} al ser una muestra pequeña
        ($\nMuestra$ observaciones), estos estadísticos son más sensibles a la
        variabilidad del muestreo que con muestras grandes; la media agrupada
        ($\mediaAgrup$) confirma la coherencia del agrupamiento.
\end{itemize}

% ====================================================================
\section{Punto 6. Parámetros poblacionales vs.\ estadísticos muestrales}
% ====================================================================

La \textbf{población} son las $\nPob$ vulnerabilidades del archivo de datos.
La muestra de $\nMuestra$ observaciones se extrajo por \textbf{muestreo
aleatorio simple sin reemplazo}. Parámetros según el formulario \FP (datos no
agrupados):
\[
\mu = \frac{1}{\nPob}\sum_{i=1}^{\nPob} x_i,
\qquad
\sigma^{2} = \frac{1}{\nPob}\sum_{i=1}^{\nPob}(x_i-\mu)^{2},
\qquad
\sigma = \sqrt{\sigma^{2}},
\qquad
CV = \frac{\sigma}{\mu}\times 100\,\text{\%}
\]

Los estadísticos muestrales $\bar{x}$, $s^{2}$, $s$ y $CV_m$ son los del
punto anterior. Para comparar: error absoluto
$E_{\text{abs}}=|\hat\theta-\theta|$ y error relativo
$E_{\text{rel}}(\%) = 100\cdot|\hat\theta-\theta|/\theta$.

\begin{table}[h!]
\centering
\caption{Comparación entre parámetros poblacionales y estadísticos muestrales
del Base Score. Fórmulas del formulario \FP.}
\medskip
\begin{tabular}{@{}lcccc@{}}
\toprule
\textbf{Estadístico} &
\makecell[c]{\textbf{Poblacional}\\($N=\nPob$)} &
\makecell[c]{\textbf{Muestral}\\($\nMuestra$ obs.)} &
\textbf{Error absol.} & \textbf{Error rel.} \\
\midrule
Media ($\mu$ / $\bar{x}$)    & $\mediaP$ & $\mediaM$ & $\errAbsMedia$ & $\errRelMedia\,\text{\%}$ \\
Varianza ($\sigma^2$ / $s^2$) & $\varP$   & $\varM$   & $\errAbsVar$   & $\errRelVar\,\text{\%}$   \\
Desv.\ estándar ($\sigma$/$s$)& $\sdP$   & $\sdM$    & $\errAbsSD$    & $\errRelSD\,\text{\%}$    \\
$CV$ (\%)                    & $\cvP$    & $\cvM$    & $\errAbsCV$    & $\errRelCV\,\text{\%}$    \\
\bottomrule
\end{tabular}
\end{table}

\subsection{Varianza de la media muestral}

El muestreo es sin reemplazo. En el formulario, el \textbf{factor de
corrección por población finita} solo aplica
\emph{``si $n/N > 0{,}05$ (muestreo sin reemplazo)''.}
Con $n=\nMuestra$ y $N=\nPob$, \CondicionFPC:
\[
\FormulaVarXbar
\]
\[
\text{Error estándar de }\bar{X} = \sqrt{\operatorname{Var}(\bar{X})} = \seXbarMedia
\]

\subsection{Interpretación}

\begin{itemize}
  \item La media muestral ($\mediaM$) aproxima a la media poblacional
        ($\mediaP$) con un error relativo de $\errRelMedia\,\text{\%}$; la
        varianza ($\errRelVar\,\text{\%}$), la desviación estándar
        ($\errRelSD\,\text{\%}$) y el $CV$ ($\errRelCV\,\text{\%}$) presentan
        errores moderados, mayores que los de muestras grandes porque
        $n=\nMuestra$ es una muestra pequeña.
  \item El error estándar teórico de $\bar X$ es $\seXbarMedia$; la diferencia
        observada $|\bar{x}-\mu|=\errAbsMedia$ equivale a $\seVsObs$ errores
        estándar, es decir, menos de un error estándar: pese al tamaño
        muestral, la muestra aleatoria simple sigue siendo representativa de
        la población.
\end{itemize}

% ====================================================================
\section{Punto 7. Cinco experimentos y sus distribuciones}
% ====================================================================

Se proponen cinco experimentos con \textbf{cinco distribuciones distintas},
todas tomadas de la tabla de distribuciones del formulario \FP. Los parámetros
se estiman con los datos del proyecto.

\begin{table}[h!]
\centering
\caption{Resumen de los cinco experimentos (fórmulas del formulario \FP).}
\medskip
\begin{tabular}{@{}clll@{}}
\toprule
\textbf{\#} & \textbf{Variable aleatoria $X$} &
\makecell[c]{\textbf{Distribución}} &
\makecell[c]{\textbf{Parámetros estimados}} \\
\midrule
1 & \makecell[l]{\# de vulnerabilidades CRÍTICAS\\entre 10 revisadas al azar}
   & Binomial            & $n=10,\ p=\expPbin$ \\
2 & \makecell[l]{\# de vulnerabilidades CRÍTICAS\\publicadas en un día}
   & Poisson             & $\lambda=\poiLam$ \\
3 & \makecell[l]{\# de CVEs que hay que revisar\\hasta encontrar el primero CRÍTICO}
   & Geométrica          & $p=\expPbin$ \\
4 & \makecell[l]{\# de ataques por red (NETWORK)\\en 20 CVEs sin reemplazo}
   & Hipergeométrica     & $\hipN,\ \hipK,\ \hipn$ \\
5 & \makecell[l]{Base Score de una vulnerabilidad\\elegida al azar}
   & Normal              & $\mu=\mediaM,\ \sigma=\sdM$ \\
\bottomrule
\end{tabular}
\end{table}

\noindent Donde $p=\expPbin$ es la proporción poblacional de severidad
\emph{CRITICAL} y, para el experimento 4, $\hipK$ es el número de
vulnerabilidades con vector de ataque \emph{NETWORK} dentro de la población de
$\hipN$ elementos.

\subsection{Experimento 1: Binomial}
\textbf{Contexto.} Un analista revisa $10$ vulnerabilidades publicadas en el
periodo, elegidas al azar. Sea $X$ el número de esas vulnerabilidades que
resultan ser de severidad \emph{CRITICAL}. Cada revisión es un ensayo de
Bernoulli con probabilidad $p=\expPbin$, independiente dada la gran diferencia
de tamaños. Según el formulario \FP:
\[
p(x) = \binom{n}{x}p^{x}(1-p)^{n-x},\qquad
X \sim \text{Bin}(n,p),\qquad
E(X)=np,\qquad \operatorname{Var}(X)=np(1-p)
\]
\[
E(X) = \binE,
\qquad
\operatorname{Var}(X) = \binV
\]
Probabilidades (con \texttt{dbinom}/\texttt{pbinom} en R):
\begin{align*}
P(X=2)     &= \binom{10}{2}(\expPbin)^{2}(1-\expPbin)^{8} = \binPdos \\
P(X\geq 1) &= 1-(1-\expPbin)^{10} = \binPgeuno
\end{align*}
Hay una probabilidad de $\binPgeuno\cdot 100\,\text{\%}$ de encontrar al menos
una vulnerabilidad crítica entre 10 revisadas.

\subsection{Experimento 2: Poisson}
\textbf{Contexto.} El periodo cubre $91$ días (01-ene-2024 a 01-abr-2024).
Sea $X$ el número de vulnerabilidades \emph{CRÍTICAS} publicadas en un día.
Se trata de un conteo de eventos en un intervalo de tiempo, el caso de uso de
la Poisson según la tabla de identificación del formulario \FP:
\[
p(x) = \frac{e^{-\lambda}\lambda^{x}}{x!},\qquad
X \sim \text{Poisson}(\lambda),\qquad
E(X) = \operatorname{Var}(X) = \lambda
\]
con $\lambda=\poiLam$ (promedio poblacional de críticas por día):
\begin{align*}
P(X=0)     &= e^{-\poiLam} = \poiPcero
   &&\text{(un día sin vulnerabilidades críticas)}\\
P(X\geq 2) &= 1-P(X=0)-P(X=1) = \poiPgeDos
   &&\text{(dos o más críticas en un día)}
\end{align*}

\subsection{Experimento 3: Geométrica}
\textbf{Contexto.} Un auditor examina vulnerabilidades una por una hasta
encontrar la primera de severidad \emph{CRITICAL}. Sea $X$ el número de
vulnerabilidades que debe revisar hasta el primer éxito; el formulario
\FP identifica esta distribución como el \emph{``número de pruebas en las que
ocurre el primer éxito''}, con probabilidad $p=\expPbin$ en cada intento:
\[
p(x) = (1-p)^{x-1}p,\qquad
X \sim \text{Geom}(p),\qquad
E(X)=\frac{1}{p},\qquad
\operatorname{Var}(X)=\frac{1-p}{p^{2}}
\]
\[
E(X) = \geoE,
\qquad
\operatorname{Var}(X) = \geoV
\]
Probabilidades (con \texttt{dgeom}/\texttt{pgeom} en R):
\begin{align*}
P(X=6)     &= (1-\expPbin)^{5}\expPbin = \geoPseis
   &&\text{(la primera crítica aparece en la 6ta.\ revisión)}\\
P(X\leq 4) &= 1-(1-\expPbin)^{3} = \geoPlecuatro
   &&\text{(encontrarla dentro de las 4 primeras revisiones)}
\end{align*}

\subsection{Experimento 4: Hipergeométrica}
\textbf{Contexto.} De la población de $\hipN$ vulnerabilidades,
$\hipK$ tienen vector de ataque \emph{NETWORK}. Un grupo selecciona
$\hipn$ vulnerabilidades \textbf{sin reemplazo}. Sea $X$ el número de esas
20 explotables por red. El formulario la define como el \emph{``número de
éxitos en una muestra $n$ tomada de una población $N$ de las cuales $k$ son
éxitos''}:
\[
p(x) = \frac{\binom{\hipK}{x}\binom{\hipN-\hipK}{\hipn-x}}{\binom{\hipN}{\hipn}},
\qquad
X \sim \text{Hipergeométrica}(\hipN,\hipK,\hipn)
\]
\[
E(X) = n\frac{K}{N} = \hipE,
\qquad
\operatorname{Var}(X)= n\frac{K}{N}\left(1-\frac{K}{N}\right)\frac{N-n}{N-1} = \hipV
\]
Probabilidades (con \texttt{dhyper}/\texttt{phyper} en R):
\begin{align*}
P(X\geq 18) &= \sum_{k=18}^{20} p(k) = \hipPgeDieciocho
   &&\text{(que al menos 18 de 20 sean por red)}\\
P(X=20) &= \frac{\binom{\hipK}{20}}{\binom{\hipN}{20}} = \hipPveinte
   &&\text{(que las 20 sean por red)}
\end{align*}

\subsection{Experimento 5: Normal}
\textbf{Contexto.} Se elige al azar una vulnerabilidad y se mide su Base
Score. Según el formulario \FP, $X\sim N(\mu,\sigma^{2})$ con densidad
\[
f(x) = \frac{1}{\sigma\sqrt{2\pi}}\,
      e^{-\frac{(x-\mu)^{2}}{2\sigma^{2}}},
\qquad E(X)=\mu,\quad \operatorname{Var}(X)=\sigma^{2}
\]
Se estiman los parámetros con la muestra: $\mu=\mediaM$, $\sigma=\sdM$.
Con la \textbf{estandarización del formulario} \FP,
\[
Z = \frac{X-\mu}{\sigma} \sim N(0,1),
\qquad
\Phi(z) = P(Z\leq z)
\]
y se resuelve con la \textbf{Tabla 4 \TP}, que entrega $P(Z>z)$ (cola
derecha); para $z<0$ se usa la simetría: $P(Z\leq -a) = P(Z>a)$.

\medskip
\textbf{1) $P(6\leq X\leq 8)$.} Estandarizando ambos límites:
\[
z_1 = \frac{6-\mediaM}{\sdM} = \normZseis,
\qquad
z_2 = \frac{8-\mediaM}{\sdM} = \normZocho
\]
\[
P(6\leq X\leq 8) = \Phi(\normZocho)-\Phi(\normZseis)
= \left[1-\TP\,P(Z>\normZocho)\right] + \TP\,P(Z>\normZseisAbs)
\]
donde se usó la simetría de la Tabla 4:
$\Phi(\normZseis)=P(Z>\normZseisAbs)$.
Consultando la Tabla 4: $P(Z>\normZocho)=\normColaZocho$ y
$P(Z>\normZseisAbs)=\normColaZseis$, de donde
\[
P(6\leq X\leq 8) = (1-\normColaZocho)+\normColaZseis = \normPseisOcho
\]

\medskip
\textbf{2) $P(X>9)$.} Estandarizando:
\[
P(X>9) = P\left(Z > \frac{9-\mediaM}{\sdM}\right)
= P(Z>\normZnueve) = \probNormMayorNueve
\]
con $P(Z>\normZnueve)=\normColaZnueve$ leído en la Tabla 4.

\medskip
\textbf{3) $P(X<5)$.} Estandarizando (nota: $z$ es negativo, se aplica la
simetría de la tabla: $P(Z<-\normZcincoAbs)=P(Z>\normZcincoAbs)$):
\[
P(X<5) = P\left(Z < \frac{5-\mediaM}{\sdM}\right)
= P(Z<\normZcinco) = P(Z>\normZcincoAbs) = \probNormMenorCinco
\]
con $P(Z>\normZcincoAbs)=\normColaZcinco$ leído en la Tabla 4.

\medskip
El $\normPseisOcho\cdot 100\,\text{\%}$ de las vulnerabilidades se espera con
puntaje entre 6 y 8 (severidad \emph{MEDIA}--\emph{HIGH}); solo
$\probNormMayorNueve\cdot 100\,\text{\%}$ superaría el 9.0 y
$\probNormMenorCinco\cdot 100\,\text{\%}$ sería inferior a 5.0.

\section*{Fuentes de fórmulas y software}
\begin{itemize}\itemsep0pt
  \item[\FP] J.~D.~Fuentes Mora, \emph{Formulario de Probabilidad y
        Estadística} (medidas de tendencia central, dispersión, momentos,
        factor de corrección por población finita, tabla de distribuciones y
        estandarización normal).
  \item[\TP] R.~E.~Walpole, \emph{Tabla 4: Áreas de curva normal}
        (probabilidad normal estándar en cola derecha $P(Z>z)$; simetría para
        $z<0$), en \emph{Introduction to Statistics}, Macmillan, 1968.
  \item Cálculos: \texttt{R/calculos/parametros.R} exporta las cifras a
        \texttt{resultados.tex}. Reproducir con:
\begin{verbatim}
cd R/calculos && Rscript parametros.R
cd ../../informe && tectonic informe.tex
\end{verbatim}
\end{itemize}

\end{document}
